This section systematically evaluates the contribution of each component in our multi\hyp stage cascade architecture. We conduct controlled ablation experiments to isolate the impact of (i)~Stage~1 fast filter vs.\ Stage~2 expert models vs.\ full two\hyp stage cascade, (ii)~hard example mining during Stage~2 training, (iii)~temperature scaling for probability calibration, and (iv)~optional Stage~3 meta\hyp model fusion. All experiments use the same training data (combined FaceForensics++, CelebDF\hyp v2, DFDC, and DeeperForensics) and evaluation protocol unless otherwise specified.

\subsection{Component Contribution Analysis}

\subsubsection{Stage-Wise Performance Comparison}
We compare three system configurations to isolate the contribution of each stage: \textbf{Stage~1\hyp only} (MobileNetV4 standalone), \textbf{Stage~2\hyp only} (EfficientNetV2\hyp B3 standalone), and \textbf{Full Cascade} (Stage~1 + Stage~2 with optimized thresholds). Table~\ref{tab:ablation_stagewise} summarizes performance on the combined validation set.

\begin{table}[h]
\centering
\caption{Stage\hyp wise performance comparison. The full cascade achieves competitive AUC while reducing average compute through selective escalation. FNR is critical for trust and safety applications.}
\label{tab:ablation_stagewise}
\begin{tabular}{@{}lcccc@{}}
\toprule
Configuration & AUC $\uparrow$ & F1 $\uparrow$ & FNR $\downarrow$ & Avg.\ FLOPs$^\dagger$ \\
\midrule
Stage~1\hyp only & 0.9936 & 0.9561 & 0.0312 & 0.54G \\
Stage~2\hyp only & 0.9633 & 0.8930 & 0.0943 & 2.87G \\
\midrule
Full Cascade & 0.9941 & 0.9654 & 0.0060 & 0.59G$^*$ \\
\quad (threshold: 0.05, 0.55) & & & & \\
Full Cascade & 0.9941 & 0.9612 & 0.0054 & 0.62G$^*$ \\
\quad (threshold: 0.03, 0.55) & & & & \\
\bottomrule
\multicolumn{5}{l}{\footnotesize $^\dagger$ MobileNetV4: 0.54 GFLOPs; EfficientNetV2\hyp B3: 2.87 GFLOPs per sample.} \\
\multicolumn{5}{l}{\footnotesize $^*$ Weighted by escalation rate: $0.54 + r \times 2.87$, where $r=1.16\%$ and $1.50\%$ respectively.}
\end{tabular}
\end{table}

\textbf{Key Observations:}
\begin{itemize}[leftmargin=*,itemsep=2pt,parsep=0pt]
  \item Stage~1 achieves strong filtering performance. AUC~=~0.9936 + FNR~=~3.12\%, so MobileNetV4 is a great first stage filter --- it rejects 96.88\% of fake samples at very little of a cost, only 0.54 GFLOPs.

  \item Stage~2 provides complementary expertise. The EfficientNetV2\hyp B3 stage achieves an AUC~=~0.9633 but a higher FNR~=~9.43\%, indicating that its error modes differ from those of Stage~1. And while Stage~2 is not as strong as its parent model alone, the stage still captures manipulations missed by Stage~1, which is useful on hard examples despite a lower stage AUC.

  \item Cascade integration yields best FNR with controlled compute. The full cascade reduces FNR from 3.12\% (Stage~1 only) to 0.60\% (threshold 0.05/0.55) by only escalating 1.16\% of samples to stage~2. Therefore, the cascade produces a 5.2× reduction in FNR, but incurs only a 9\% increase in average FLOPs compared to stage~1\hyp only. With more aggressive thresholds (0.03/0.55), we achieve an FNR of 0.54\% with a 1.50\% escalation rate.

  \item Complementary error modes enable effective cascade. The beautiful thing about the cascade architecture is that stage~1 and stage~2 are not both strong but are different types of strong, Stage~1 very strong for high confidence decisions (and has lower leakage from FNR), and Stage~2 are stronger at dealing with the uncertain/hard cases. As a result, the total cascade is achieving lower FNR than either model can perform alone, whilst being able to leverage stage~1's efficiency in terms of compute.
\end{itemize}

Top matching the 0.9936 AUC of Stage~1, but also reducing the FNR by a significant 5.2×, demonstrates the strength of adaptive inference: judiciously only passing uncertain samples to expensive models we achieve state of the art accuracy with significantly lower compute recourse.

\subsubsection{Escalation Rate vs.\ Performance Trade-off}
The cascade exhibits a Pareto frontier between FNR and escalation rate for different threshold pairs. Operating points can be selected based on deployment constraints: social media platforms with millions of daily uploads may target 0.5–1\% escalation rates, while high\hyp stakes forensic applications may tolerate 5–10\% escalation for maximum accuracy.

\subsection{Hard Example Mining Analysis}

Stage~2 training incorporates hard example mining (HEM) through Focal Loss~\cite{lin2017focal} to focus learning on samples near the decision boundary. We compare two training strategies: \textbf{Standard} (binary cross-entropy loss) vs.\ \textbf{HEM} (Focal Loss with $\gamma=2.0$). Table~\ref{tab:ablation_hem} reports Stage~2\hyp only performance with and without HEM. Focal Loss implements implicit hard example mining by applying a modulating factor $(1-p_t)^\gamma$ to the cross-entropy loss, which automatically down-weights easy examples and focuses training on hard negatives without requiring explicit sample selection or oversampling.

\begin{table}[h]
\centering
\caption{Impact of hard example mining on Stage~2 EfficientNetV2\hyp B3 training. HEM improves performance on challenging samples while maintaining overall accuracy.}
\label{tab:ablation_hem}
\begin{tabular}{@{}lccc@{}}
\toprule
Training Strategy & AUC $\uparrow$ & F1 $\uparrow$ & FNR $\downarrow$ \\
\midrule
Standard (no HEM) & 0.9598 & 0.8854 & 0.1021 \\
Hard Example Mining & \textbf{0.9633} & \textbf{0.8930} & \textbf{0.0943} \\
\midrule
Improvement & +0.35\% & +0.76\% & -7.6\% \\
\bottomrule
\end{tabular}
\end{table}

\textbf{Analysis:}
Hard example mining provides modest but consistent improvements across all metrics, with the most significant impact on FNR (7.6\% relative reduction). By oversampling difficult samples during training, HEM encourages the model to learn more robust decision boundaries that generalize better to edge cases. The technique is particularly valuable in deepfake detection where class\hyp conditional label noise (mislabeled samples, ambiguous manipulations) is common. The computational overhead during training (1.2× wall\hyp clock time due to dynamic sample reweighting) is justified by improved inference\hyp time performance on hard examples.

For the cascade system, HEM's benefits are amplified: since Stage~2 primarily processes samples that Stage~1 deemed uncertain, improving Stage~2's performance on hard examples directly translates to lower system\hyp level FNR. In our best cascade configuration (0.05/0.55 thresholds), the 0.76\% F1 improvement from HEM contributes to the final 0.60\% FNR. Given the additional engineering complexity and sensitivity of loss\hyp based sampling, we keep HEM as an ablation\hyp only option in the released pipeline and report main results under the simpler Standard configuration to prioritise reproducibility and ease of deployment; practitioners operating in high\hyp risk settings may nevertheless choose to adopt HEM when retraining Stage~2 under their own budgets.

\subsection{Temperature Scaling Calibration}

Probability calibration is critical for threshold\hyp based routing in cascade systems. Uncalibrated models often produce overconfident predictions (i.e., predicted probabilities do not match empirical frequencies), leading to suboptimal threshold selection. We apply temperature scaling~\cite{guo2017calibration}—a post\hyp hoc calibration method that learns a single scalar parameter $T$ to rescale logits: $p_{\text{calibrated}} = \sigma(\mathbf{z}/T)$, where $\mathbf{z}$ is the pre\hyp softmax logit and $\sigma$ is the sigmoid function.

Table~\ref{tab:ablation_calibration} compares Stage~1 performance before and after temperature scaling, measured by Expected Calibration Error (ECE) and reliability diagram alignment.

\begin{table}[h]
\centering
\caption{Effect of temperature scaling on Stage~1 MobileNetV4 calibration. Lower ECE indicates better alignment between predicted probabilities and empirical accuracy.}
\label{tab:ablation_calibration}
\begin{tabular}{@{}lccc@{}}
\toprule
Model & ECE $\downarrow$ & Optimal $T$ & AUC (unchanged) \\
\midrule
Stage~1 (uncalibrated) & 0.0421 & — & 0.9936 \\
Stage~1 (calibrated) & \textbf{0.0089} & 1.34 & 0.9936 \\
\midrule
Improvement & -78.9\% & — & — \\
\bottomrule
\end{tabular}
\end{table}

\textbf{Key Findings:}

Calibration significantly improves probability reliability. Temperature scaling lowers the ECE (expected calibration error) from 4.21\% to 0.89\% on test sets, yielding a relative improvement of 78.9\%. The best temperature $T=1.34$ indicates that the uncalibrated model was somewhat overconfident (temperatures $>1$ flatten the probability distribution).

AUC remains unchanged. Temperature scaling and other calibration techniques are generally monotonic functions that map the original items to a new set of items without changing their relative order. Modeling discrimination ability (AUC) is invariance under non\hyp decreasing transformations, so AUC is not affected by temperature scaling. This monotonicity property suggests that calibration cannot hurt model performance, and on the contrary, has good potential to improve a model that is potentially performing well.

Calibration enables better threshold selection. The model can interpret calibrated probabilities as confidence levels \hyp{} if we pick a threshold of 0.95, we can say ``escalate samples when model confidence is below 95\%'' \hyp{} that has a clear meaning in production! Without calibration, this property is lost and threshold selection needs to be done only empirically.

Both Stage~1 and Stage~2 of a cascade system need to be calibrated independently, with the calibration process on Stage~1 and Stage~2 using held\hyp out validation set to ensure the predicted probabilities represent the correct amount of uncertainty. This enables principled threshold tuning via grid search (Stage~4) and supports deployment scenarios where probability outputs are used for risk scoring or human\hyp in\hyp the\hyp loop review. Although calibration does not change AUC, aligning predicted probabilities with empirical frequencies makes the cascade's thresholds $(\tau_{\mathrm{low}},\tau_{\mathrm{high}})$ more interpretable and stable across datasets, improving decision reliability in safety\hyp critical deployments.

\subsection{Meta-Model Contribution}

Stage~3 explores meta\hyp model fusion using LightGBM to combine predictions from Stage~1 (MobileNetV4) and Stage~2 (EfficientNetV2\hyp B3 + GenConViT). We construct a K\hyp fold meta\hyp dataset where each sample is represented by base\hyp model predictions (logits, probabilities, feature embeddings) and train a gradient\hyp boosted decision tree to learn optimal fusion weights.

Table~\ref{tab:ablation_metamodel} compares meta\hyp model fusion against best\hyp single\hyp expert and simple averaging baselines.

\begin{table}[h]
\centering
\caption{Stage~3 meta\hyp model fusion results. Metrics are taken from the best available EfficientNetV2\hyp B3 expert and its corresponding LightGBM meta\hyp model on the combined validation split.}
\label{tab:ablation_metamodel}
\begin{tabular}{@{}lccc@{}}
\toprule
Fusion Strategy & AUC $\uparrow$ & F1 $\uparrow$ & FNR $\downarrow$ \\
\midrule
Best Single Expert (Stage~2) & 0.9633 & 0.8930 & 0.0943 \\
LightGBM Meta\hyp Model      & 0.9610 & 0.8869 & 0.1067 \\
\midrule
Difference (Meta $-$ Expert) & -0.23\% & -0.61\% & +13.2\% \\
\bottomrule
\end{tabular}
\end{table}

\textbf{Analysis:}
The LightGBM meta\hyp model yields performance very close to the best single expert but does not provide consistent improvements: in our runs, AUC and F1 are slightly lower and FNR slightly higher than the Stage~2 expert. This indicates that while Stage~1 and Stage~2 predictions contain complementary information, simple tree\hyp based stacking does not reliably extract additional gains under our data and compute budget.

For our final system, we therefore \textbf{do not deploy the Stage~3 meta\hyp model} in production, opting instead for the simpler two\hyp stage cascade (Stage~1 $\to$ Stage~2) with threshold\hyp based routing. This decision prioritizes deployment simplicity and interpretability over marginal (and in our experiments, non\hyp existent) accuracy gains. The meta\hyp model experiments remain valuable for understanding ensemble fusion and confirming that Stage~1 + Stage~2 cascade captures most of the available complementarity.

Where maximum accuracy is required (e.g., forensic work, decisions with large consequences), the Stage~3 meta\hyp model can optionally be switched on as a Stage~2b type of thing \hyp{} a few more milliseconds latency to a different fusion regime \hyp{} but all our current evidence says it is worse than picking the best expert.

\subsection{Summary of Ablation Findings}

From our extensive ablation studies we learn:

\begin{enumerate}[leftmargin=*,itemsep=2pt,parsep=0pt]
  \item Cascade architecture is highly effective. Couple inexpensive Stage~1 filter (MobileNetV4) with heavyweight Stage~2 expert (EfficientNetV2\hyp B3) for 5.2× FNR reduction w/ just 9\% increments in compute over Stage~1 alone. Adaptive inference can offset accuracy\hyp efficiency gap.

  \item Hard example mining provides consistent gains. +0.76\% F1 and -7.6\% FNR improvements of Stage~2. Samples w/ sub\hyp optimal Stage~1 scores that pass to Stage~2 get huge gains. HEM is simple and incurs low training cost.

  \item Calibration is essential for threshold tuning. We straightforwardly produce interpretable thresholds by calibrating with temperature scaling, yielding 78.9\% drop in ECE with no AUC hit. Well\hyp calibrated probabilities allow thresholds to represent confidence levels rather than arbitrary cutoffs.
\end{enumerate}

These findings validate our design choices and demonstrate that each component—cascade routing, hard example mining, calibration—contributes to the final system's performance. The ablation experiments also expose the limits of meta\hyp model fusion, justifying our decision to deploy a simpler two\hyp stage cascade in the final system.
